\chapter{IMPLEMENTASI DAN SIMULASI}

\section{Struktur Program dan Dependensi}
Program menggunakan Go 1.23 atau lebih baru. Cipher dan mode operasi
diimplementasikan dalam repositori tanpa pustaka kriptografi pihak ketiga.
Pustaka standar \texttt{crypto/rand} dipakai oleh CLI untuk membangkitkan
IV; \texttt{math/rand} hanya dipakai untuk eksperimen yang dapat diulang.

\begin{table}[htbp]
  \centering\small
  \caption{Komponen kode sumber.}
  \begin{tabularx}{\linewidth}{@{}p{0.48\linewidth}Y@{}}
    \toprule Lokasi & Peran\\ \midrule
    \path{src/cipher/cipher.go} & Kontrak Block, Factory, dan Tracer.\\
    \path{src/cipher/brainrot/} & Cipher, transformasi, tabel, dan ekspansi kunci.\\
    \path{src/modes/modes.go} & Lima mode, PKCS\#7, dan pencacah CTR.\\
    \path{src/cmd/crypt/main.go} & CLI enkripsi teks dan dekripsi hex.\\
    \path{src/cmd/analyze/main.go} & CLI analisis berkas uji dan keluaran JSON.\\
    \path{src/analysis/} & Statistik distribusi dan eksperimen perubahan bit.\\
    \path{src/tests/} & Pengujian infrastruktur dengan cipher XOR uji.\\
    \bottomrule
  \end{tabularx}
\end{table}

\section{Antarmuka dan Format Data}
CLI \texttt{crypt} menerima \texttt{-op}, \texttt{-mode},
\texttt{-key-hex}, dan \texttt{-input}. Kunci harus berupa 32 karakter
heksadesimal. Saat enkripsi, masukan berupa teks; saat dekripsi,
masukan berupa hex hasil enkripsi. Mode selain ECB menyimpan IV 16 byte
di depan badan cipherteks: $\mathrm{IV}\Vert C$. Mode ECB hanya menyimpan $C$.
Mode dan kunci tidak disertakan dalam keluaran; pengguna memasukkan
keduanya lagi saat dekripsi.

\begin{lstlisting}[language={},numbers=none,caption={Contoh simulasi CBC dari root repositori pada Bash.}]
key_hex=00112233445566778899aabbccddeeff
ciphertext=$(go -C src run ./cmd/crypt -op encrypt \
  -mode CBC -key-hex "$key_hex" -input "Halo Brainrot")
go -C src run ./cmd/crypt -op decrypt \
  -mode CBC -key-hex "$key_hex" -input "$ciphertext"
\end{lstlisting}

\Isi{Cantumkan keluaran aktual atau tangkapan terminal. Ulangi simulasi
untuk ECB, CBC, CFB, OFB, dan CTR. Deskripsikan antarmuka akhir jika berubah.}

\section{Mode Operasi, IV, dan Counter}
\begin{table}[htbp]
  \centering\small
  \caption{Perilaku mode operasi dalam kode saat ini.}
  \label{tab:mode}
  \begin{tabularx}{\linewidth}{@{}lp{0.23\linewidth}Y@{}}
    \toprule Mode & IV / counter & Blok akhir dan umpan balik\\ \midrule
    ECB & Tidak digunakan & Blok independen; PKCS\#7.\\
    CBC & IV 16 byte & XOR dengan blok cipherteks sebelumnya; PKCS\#7.\\
    CFB & IV 16 byte & Umpan balik blok penuh (CFB-128); potongan akhir
      memakai sebagian keystream tanpa padding.\\
    OFB & IV 16 byte & Umpan balik keluaran cipher; tanpa padding.\\
    CTR & Counter awal 16 byte & Pencacah big-endian 128 bit;
      menolak overflow jika masih ada data; tanpa padding.\\
    \bottomrule
  \end{tabularx}
\end{table}
Untuk CBC, CFB, dan OFB, IV dibuat oleh CLI menggunakan
\texttt{crypto/rand}. Nilai yang sama menjadi counter awal CTR.
Dekripsi CFB, OFB, dan CTR menggunakan fungsi enkripsi blok untuk
membangkitkan keystream.

\Isi{Jelaskan alasan pengelolaan IV/counter dan mekanisme penyimpanannya
pada format berkas akhir. Catat cara mencegah penggunaan ulang keystream.}

\section{Penanganan Padding}
Untuk ECB dan CBC dengan $b=16$, panjang badan cipherteks adalah
$16(\lfloor L/16\rfloor+1)$ byte. Untuk CFB, OFB, dan CTR, panjang badan
cipherteks sama dengan panjang plainteks. Panjang IV tidak termasuk dalam
kedua ukuran tersebut. Dekripsi ECB/CBC memeriksa panjang cipherteks,
nilai byte padding terakhir, dan kesamaan seluruh byte padding.

\Isi{Justifikasi PKCS\#7 dan lengkapi bukti untuk masukan 0, 1, 15, 16,
17, dan 32 byte. Catat panjang plainteks, padding, badan cipherteks,
dan hasil pemulihan; sertakan kasus padding tidak valid.}

\section{Simulasi Berkas Teks dan Biner}
Fungsi mode menerima slice byte, sedangkan CLI analisis membaca berkas.
CLI \texttt{crypt} saat ini hanya menerima teks dari argumen dan belum
menyediakan pembacaan serta penulisan berkas biner sembarang.

\Isi{Lengkapi bagian ini setelah alur enkripsi/dekripsi berkas tersedia.
Gunakan berkas teks dan biner pada seluruh mode; jangan menganggap
simulasi argumen teks sebagai bukti dukungan berkas biner.}

\begin{table}[htbp]
  \centering\small
  \caption{Isian bukti pemulihan data untuk seluruh mode.}
  \label{tab:simulasi}
  \begin{tabularx}{\linewidth}{@{}lYYYY@{}}
    \toprule Mode & Teks: berkas / ukuran & Biner: berkas / ukuran &
      Panjang keluaran & Byte hasil dekripsi sama\\ \midrule
    ECB & --- & --- & --- & ---\\
    CBC & --- & --- & --- & ---\\
    CFB & --- & --- & --- & ---\\
    OFB & --- & --- & --- & ---\\
    CTR & --- & --- & --- & ---\\
    \bottomrule
  \end{tabularx}
\end{table}

\section{Verifikasi Integritas dan Kasus Gagal}
Analisis menggunakan \texttt{bytes.Equal} untuk memeriksa bahwa hasil
dekripsi sama persis dengan masukan. Kode saat ini belum menyimpan
tag autentikasi atau mekanisme verifikasi perubahan cipherteks. Pemeriksaan
padding hanya memvalidasi bentuk padding, sehingga tidak mencakup
seluruh perubahan pada isi pesan.

\Isi{Jelaskan mekanisme integritas yang dipilih pada implementasi akhir
sesuai spesifikasi. Bedakan bukti round-trip dengan deteksi modifikasi
data. Sertakan hasil pengujian kunci salah, data berubah, IV berubah,
data terpotong, dan masukan tidak valid.}
